\section{Pre-registration protocol}\label{sec:prereg}

The project keeps a dated, append-only decision log. A hypothesis, target, lag, sign
or acceptance bar is committed the day its entry lands, before the corresponding fit;
a later change is a visibly dated superseding entry that names the result which
prompted it. The log is the instrument that turns a null into a finding: without it,
a reader cannot distinguish ``we tested and found nothing'' from ``we kept looking
until nothing crossed a line, and reported the line''. Table~\ref{tab:prereg} lists
the entries that fixed each model below.

\begin{table}[t]
\centering\small
\caption{Pre-registration entries. The commit column is the honest part: see the text.}
\label{tab:prereg}
\begin{tabular}{llp{7.2cm}l}
\toprule
entry & dated & what was fixed before the fit & commit \\
\midrule
  D-001--D-003 & 2026-08-10 & A1 target, kernel, replay protocol & \texttt{4239d13} \\
  D-016 & 2026-08-12 & A2 design lock & \texttt{4239d13} \\
  D-019 & 2026-08-12 & A5 design lock & \texttt{4239d13} \\
  D-022 & 2026-08-12 & A4 design lock & \texttt{4239d13} \\
  D-025 & 2026-08-12 & H1 horizon test pre-specified & \texttt{4239d13} \\
  D-028 & 2026-09-05 & Part B null + FWL: target, controls, HAC lag, signs & \texttt{e483a1b} \\
  D-031 & 2026-09-05 & H2/H3/H4: signs, Bonferroni bar, holdout & \texttt{c489550} \\
\bottomrule
\end{tabular}

\end{table}

Two consequences of the protocol are visible in the results. A pre-registered sign
falsified by the data is reported as falsified (Section~\ref{sec:parta}); and a test
that produced an apparent 87\% improvement over its null was withdrawn, because the
protocol's own comparability condition failed (Section~\ref{sec:parta}). Neither
would have happened without a committed criterion to fail against.

\paragraph{What the log does not establish.} The log is an ordinary text file in the
project's version control, and its git history does not independently prove the
ordering it asserts. The entries for the physical nowcasts (D-001 through D-026) were
committed together, weeks after the work they describe; the two spread
pre-registrations (D-028 and D-031) were each committed in the same change as their
result. The \emph{substance} of the pre-registration --- the signs, bars, horizon and
holdout fixed in each entry's text --- was written before the corresponding run, but
a sceptical reader has only the author's word for it. We state this rather than have
it discovered, and note that a third-party timestamped registry would have cost
nothing. Future entries are committed and pushed before the run they govern.
